\section{Space Complexity --- Machines --- ARMRun}
%%% generated by the blueprint pipeline from TCSlib.Complexity.SpaceComplexity.Machines.ARMRun
\subsection*{Space Complexity --- Machines --- ARMRun: Introduction}

Declarations of the Lean module \texttt{TCSlib.Complexity.SpaceComplexity.Machines.ARMRun}.

\subsection*{Space Complexity --- Machines --- ARMRun: Results}

\begin{definition}[Machine precondition]
\lean{Complexity.LogProg.Pre}
\label{Complexity.LogProg.Pre}
\leanok
\uses{Complexity.LogProg.AConf, Complexity.LogProg.ARM, Complexity.LogProg.CleanRun, Complexity.LogProg.ValidPair, Complexity.LogProg.ValidPlain, Complexity.LogProg.callSegs, Complexity.LogProg.vword, Turing.MultiTapeTM}
Let $A$ be a program with $m$ registers, an oracle, a Boolean multi-tape Turing machine
$D$, a choice of initial state for each decider, and a natural-number bound $B$. A
configuration satisfies the precondition if, at a register-comparison instruction, the
compared registers are distinct; at an input-payload comparison, the input has the
corresponding valid encoding; and at a call instruction, the argument registers are
distinct and $D$, started in the assigned state on the call’s encoded virtual input,
halts with the oracle’s answer, leaves all work tapes blank with all work-tape heads at
zero, and keeps every work-tape head within distance $B$ of zero throughout the run,
including its halting time. Other configurations satisfy the precondition automatically.
\end{definition}

\begin{lemma}[Successor step of an abstract run]
\lean{Complexity.LogProg.arun_succ}
\label{Complexity.LogProg.arun_succ}
\leanok
\uses{Complexity.LogProg.AConf, Complexity.LogProg.ARM, Complexity.LogProg.arun, Complexity.LogProg.astep}
\difficulty{1}
For any register program $A$, oracle function $o$, input $x$, configuration $a$, and
natural number $n$, running for $n+1$ steps from $a$ gives the same configuration as
first taking one abstract step from $a$ and then running for $n$ more steps.
\end{lemma}

\begin{lemma}[Halted configurations remain fixed]
\lean{Complexity.LogProg.arun_halted}
\label{Complexity.LogProg.arun_halted}
\leanok
\uses{Complexity.LogProg.ARM, Complexity.LogProg.arun, Complexity.LogProg.arun_succ}
\difficulty{4}
For any natural numbers $m,d$, type of labels $\Lambda$, program $A$ assigning an
instruction to each label, oracle $o : \{0,\ldots,d-1\} \to \mathrm{List}(\{0,1\}) \to
\{0,1\}$, input bit string $x$, register assignment $v : \{0,\ldots,m-1\} \to
\mathbb{N}$, stored answer $res \in \mathrm{Option}(\{0,1\})$, and number of steps $n$,
iterating the program’s transition function $n$ times from the configuration with no
active label, registers $v$, and answer $res$ leaves that configuration unchanged.
\end{lemma}

\begin{definition}[Simulation invariant]
\lean{Complexity.LogProg.Inv}
\label{Complexity.LogProg.Inv}
\leanok
\uses{Complexity.LogProg.ARM, Complexity.LogProg.CallOK, Complexity.LogProg.Ph, Complexity.LogProg.armProg, Turing.Cfg, Turing.MultiTapeTM}
Let $A$ be a program with $m$ registers, let $l_0$ be a control label, let an oracle
assign Boolean answers to deciders on bit strings, and let $D$ be a Boolean multi-tape
Turing machine with an initial state assigned to each decider. For natural numbers $B$
and $W$, a configuration of the compiled machine on input $x$ satisfies the invariant if
every work-tape head position lies between $-1$ and $W$, inclusive, and the
configuration satisfies the call-validity condition for the program compiled from $A$
starting at $l_0$, with head-position bounds from $-1$ through $W$ and call bound $B$.
\end{definition}

\begin{lemma}[Invariant at call-free bounded configurations]
\lean{Complexity.LogProg.inv_of_mid}
\label{Complexity.LogProg.inv_of_mid}
\leanok
\uses{Complexity.LogProg.ARM, Complexity.LogProg.Inv, Complexity.LogProg.Mid, Complexity.LogProg.Ph, Complexity.LogProg.armProg, Turing.Cfg, Turing.MultiTapeTM}
\difficulty{3}
Let a register program, its initial label, an oracle, a decider machine, and a
starting-state assignment be fixed, together with natural numbers $B,W$ and a
configuration $c$ on a Boolean input. If no active state of $c$ is an oracle-call state
and every register head position lies between $-1$ and $W$, then every register head
position lies between $-1$ and $W$, and each active oracle call satisfies the following
conditions: its argument registers are distinct, the input head is at position $1$, and
each argument register stores its represented word contiguously from cell zero with its
head at zero, with the word length bounded below by $1$ and above by $W$. Moreover, the
decider selected by that call, started on the virtual input assembled from the call
specification and the represented register words, halts with the oracle's answer as its
sole output, leaves all work tapes blank and their heads at zero, and keeps every
work-tape head within distance $B$ of zero throughout the run.
\end{lemma}

\begin{lemma}[Preservation of the invariant through a halting run]
\lean{Complexity.LogProg.inv_append}
\label{Complexity.LogProg.inv_append}
\leanok
\uses{Complexity.LogProg.ARM, Complexity.LogProg.Inv, Complexity.LogProg.Ph, Complexity.LogProg.SimTo, Complexity.LogProg.armProg, Complexity.LogProg.inv_of_mid, Complexity.LogProg.rrun, Complexity.LogProg.rrun_add, Turing.Cfg, Turing.MultiTapeTM}
\difficulty{5}
Let a register program, its initial label, an oracle, a decider machine, and a
starting-state assignment be fixed, together with natural numbers $B,W$, two
configurations $c,c'$ on the same Boolean input, an output word $o$, and a bit $b$.
Suppose the program simulates $c'$ from $c$ in finitely many steps, with no oracle-call
state and all register heads between $-1$ and $W$ at every intermediate step. Also
suppose that from $c'$ it reaches a halting configuration whose output is $o$ followed
by $b$, whose register heads lie between $-1$ and $W$, and whose every earlier
configuration satisfies the invariant. Then from $c$ it reaches a halting configuration
with that same output and head bounds, and every earlier configuration satisfies the
invariant.
\end{lemma}

\begin{theorem}[Simulation of a halting register computation]
\lean{Complexity.LogProg.arm_run}
\label{Complexity.LogProg.arm_run}
\leanok
\uses{Complexity.LogProg.ARM, Complexity.LogProg.CleanRun, Complexity.LogProg.HaltsWith, Complexity.LogProg.Inv, Complexity.LogProg.Pre, Complexity.LogProg.SimTo, Complexity.LogProg.ValidPair, Complexity.LogProg.ValidPlain, Complexity.LogProg.armProg, Complexity.LogProg.armProg_call, Complexity.LogProg.arun, Complexity.LogProg.arun_halted, Complexity.LogProg.arun_succ, Complexity.LogProg.aseam, Complexity.LogProg.astep, Complexity.LogProg.callSegs, Complexity.LogProg.insCall, Complexity.LogProg.inv_append, Complexity.LogProg.inv_of_mid, Complexity.LogProg.length_bits_mono, Complexity.LogProg.regWords, Complexity.LogProg.rrun, Complexity.LogProg.rrun_add, Complexity.LogProg.rrun_zero, Complexity.LogProg.sim_call, Complexity.LogProg.sim_clr, Complexity.LogProg.sim_dec, Complexity.LogProg.sim_half, Complexity.LogProg.sim_inc, Complexity.LogProg.sim_jeq, Complexity.LogProg.sim_jeqFst, Complexity.LogProg.sim_jeqIn, Complexity.LogProg.sim_jeqSnd, Complexity.LogProg.sim_jodd, Complexity.LogProg.sim_jz, Complexity.LogProg.sim_ret, Complexity.LogProg.sim_valP, Complexity.LogProg.sim_valQ, Complexity.LogProg.tapeWord_bufferTape, Complexity.LogProg.vword, Turing.MultiTapeTM}
\difficulty{7}
Fix a Boolean input $x$, an arithmetic register program $A$, its initial state $l_0$, an
oracle, a multi-tape machine $D$, a choice of initial state for $D$ for each oracle
index, and bounds $B,W\in\mathbb{N}$. Suppose that an $N$-step register computation from
state $l$ and valuation $v$, with no answer yet, has halted with answer $b$. Assume that
at every earlier step its configuration satisfies the register-program precondition for
$D$ and the chosen oracle, and each register value plus one has binary length at most
$W$. Then the tape machine compiled from $A$ can run from the initial tape configuration
encoding $x,l,v$ and any output word $out$ to a halting configuration whose output is
$out$ followed by $b$, with every work-tape head position between $-1$ and $W$
inclusive; its invariant holds at every earlier step.
\end{theorem}

\begin{lemma}[Initial configuration as a simulation seam]
\lean{Complexity.LogProg.init_eq_aseam}
\label{Complexity.LogProg.init_eq_aseam}
\leanok
\uses{Complexity.LogProg.aseam, Turing.Cfg.init}
\difficulty{3}
For any starting state $l_0$, the initial machine configuration on input $x$ with state
$(l_0,\mathrm{start})$ equals the seam configuration on $x$ with base state $l_0$, all
registers zero, and empty output.
\end{lemma}

\begin{theorem}[Space bound for compiled register programs]
\lean{Complexity.LogProg.arm_space}
\label{Complexity.LogProg.arm_space}
\leanok
\uses{Complexity.LogProg.ARM, Complexity.LogProg.Pre, Complexity.LogProg.armProg, Complexity.LogProg.arm_run, Complexity.LogProg.arun, Complexity.LogProg.compileFinTM, Complexity.LogProg.compile_space, Complexity.LogProg.init_eq_aseam, Turing.FinTM.ComputesInTime, Turing.MultiTapeTM, Turing.MultiTapeTM.initCfg, Turing.MultiTapeTM.spaceUsed}
\difficulty{5}
Let $A$ be a register program with $m$ registers and access to $d$ Boolean oracles, let
$x$ be a Boolean input, and let $D$ be a $k_D$-work-tape machine with designated initial
states $q_0(j)$ for each oracle. Suppose the run of $A$ from its initial state, zeroed
registers, and no stored answer halts after $N$ steps with answer $b$. Suppose also that
at every time $t<N$, its configuration satisfies the precondition for $A$, and every
register value $v(r)$ satisfies $\lvert\mathrm{bits}(v(r)+1)\rvert\le W$. Then the
compiled machine halts on $x$ with output $b$ at some time, and its total space use is
at most $m(W+2)+k_D(2B+1)$.
\end{theorem}
